\section{Gap in hopping units}
\label{sec:h-mass}

The normalised uniform mode of \(8\)-regular \(J_2\) sits at eigenvalue \(-8\),
independent of the number of vertices.
Define \(m=-8-E_0\) when \(E_0<-8\) and \(m=0\) otherwise.
Proposition~\ref{prop:clique} gives \(m\ge n-9\) for an embedded clique of
order \(n\ge 10\).
The trial vector uses only the edges inside its support.
Extra edges from that support to the bath do not raise the Rayleigh quotient.
Removing edges lowers the all-ones quotient toward the interior of the band,
so the same trial vector does not produce a gap from under-coordination.
This \(m\) is an eigenvalue difference, not an extensive energy:
\(E_\psi=\langle\psi|H|\psi\rangle\) for a normalised state remains \(O(1)\).
A conversion into a rest energy in the units of I1 is a further identification.
The dimensionless difference \(E_0+8\) does not need it.
The same gap is what Section~\ref{sec:h-merge} distinguishes from \(\sum R\).
